\makeatletter

\newcommand{\InitialData}[1]{\def\@InitialData{#1}}
\newcommand{\ReportContents}[1]{\def\@ReportContents{#1}}
\newcommand{\ExpectedPages}[1]{\def\@ExpectedPages{#1}}

% ==> Переоформление титульного листа
\renewcommand{\maketitle}{
    \begin{titlepage}
        \begin{center}
            МИНОБРНАУКИ РОССИИ \\
            \rule[0.5ex]{0.6\textwidth}{1.2pt} \\
            СПбГЭТУ «ЛЭТИ» ИМ. В.И. УЛЬЯНОВА (ЛЕНИНА) \\
            \rule[0.5ex]{0.6\textwidth}{1.2pt} \\
            Кафедра \@Department
        \end{center}

        \vfill

        \begin{center}
            \textbf{
                \MakeUppercase{Курсовая работа} \\
                По дисциплине «\@Discipline» \\
                Тема: \MakeUppercase{\@WorkTitle} \\
            }
        \end{center}

        \vfill

        \noindent
        \begin{tabularx}{\textwidth}{l X r}
            Студент гр. \@Group \ifdefempty{\@Variant}{}{, вар. \@Variant}& \hrulefill & \@StudentName \\\\
            Преподаватель    & \hrulefill & \@TeacherName
        \end{tabularx}

        \vfill

        \begin{center}
            Санкт-Петербург \\
            \@Year \\
        \end{center}
        \thispagestyle{empty}
    \end{titlepage}
    \newpage
    \setcounter{page}{2}
}
% <==

% ==> Оформление страницы задания
\newcommand{\makeassignment}{
    \begin{center}
        \textbf{ЗАДАНИЕ \\ НА КУРСОВУЮ РАБОТУ}
    \end{center}

    % \vspace{1em}

    \noindent Студент \@StudentName \\
    
    \noindent Группа \@Group \\
    
    \noindent Тема работы: \@WorkTitle \\

    % \vspace{1em}

    \noindent Исходные данные: \\
    \@InitialData \\

    % \vspace{1em}

    \noindent Содержание пояснительной записки: \\
    \@ReportContents \\

    % \vfill

    \noindent Предполагаемый объем пояснительной записки: \\
    Не менее 20~страниц. \\

    \noindent \makebox[5cm][l]{Дата выдачи задания:}  \rule{5cm}{0.4pt} \\
    \noindent \makebox[5cm][l]{Дата сдачи реферата:}  \rule{5cm}{0.4pt} \\
    \noindent \makebox[5cm][l]{Дата защиты реферата:} \rule{5cm}{0.4pt}

    % \vspace{2em}
    \vfill

    \noindent
    \begin{tabularx}{\textwidth}{l X r}
        Студент    & \hrulefill & \@StudentName \\\\
        Преподаватель  & \hrulefill & \@TeacherName
    \end{tabularx}

    \thispagestyle{plain} % или empty, если не нужна нумерация на этой странице
    \newpage
}
% <==

\makeatother